\section{The chamber regularization}\label{sec:regularization}

The photon field is a real one-cochain on the cubic lattice. The fermion
has one site in each chamber of the lattice mirrors. This choice has a
specific consequence: a fermion bond crossing a mirror joins a site to
its reflected site and is perpendicular to the mirror. The hopping term
then admits a positive reflected expansion. We prove this expansion and
the resulting finite-volume reflection positivity below. The construction
also requires a nonzero partition function when the bare mass has
undetermined sign. For this purpose we identify a nonzero coefficient
of the fermion determinant, independent of the mass. We reproduce the
chamber construction and this normalization argument from
\cite[Sections~2.1--2.3, 3.1--3.5 and 4.1--4.4]{ChamberQED}.

Throughout this section, $0<\rho<1$, $q>0$, and the mass parameter $M$
is real. The hopping weights $w_0,\ldots,w_3$ are positive. All constants may depend on these fixed parameters unless a uniform
bound is stated.

\subsection{The photon field and the flux sum}\label{subsec:reg-photon}

We orient each link $(x,\mu)$ from $x$ to $x+e_\mu$ and each plaquette
$(x;\mu\nu)$, $\mu<\nu$, by the ordered pair $(e_\mu,e_\nu)$.
The real $k$-cochains are denoted by $C^k$, and the integer cochains by
$C^k(\Z)$. The coboundary $d$ and its adjoint $d^*$ use these
orientations. Thus
\begin{equation}\label{eq:reg-coboundary}
 (dA)_{\mu\nu}(x)=A_\mu(x)+A_\nu(x+e_\mu)
                     -A_\mu(x+e_\nu)-A_\nu(x).
\end{equation}
An oriented link may also be written as a pair of its endpoints;
$A(y,x)=-A(x,y)$ and $A(x,x)=0$. Pullback by a lattice isometry includes
the orientation of the image cell. In particular, pullback commutes
with $d$. For a coordinate reflection, this gives a minus sign on
the curvature components containing the normal direction and a
plus sign on the tangential component. These are the two-form
reflection signs used throughout.

Let $\Gamma\subset\Z^3$ be a rank-three subgroup with
\[
 \min\{|g|_\infty:g\in\Gamma\setminus\{0\}\}\ge4.
\]
We write $C^k(\Gamma)$ for the cochains on $\Z^3/\Gamma$ and identify
them with their periodic extensions. The torus has
$V=[\Z^3:\Gamma]$ photon sites, $3V$ links and $3V$ plaquettes. The
inner product on real cochains is the sum over positively oriented
cells, and $\|\cdot\|$ denotes its norm.

A photon configuration is a class in
\begin{equation}\label{eq:reg-configuration}
 \operatorname{Conf}(\Gamma)
 =\{(A,z):A\in C^1(\Gamma),\ z\in C^2(\Gamma;\Z),\ dz=0\}/\!\sim,
 \qquad
 (A,z)\sim\left(A+\frac{2\pi}{q}k,z-dk\right),
\end{equation}
where $k\in C^1(\Gamma;\Z)$. Its curvature and link phases are
\begin{equation}\label{eq:reg-curvature}
 F=dA+\frac{2\pi}{q}z,
 \qquad u_b=\exp(iqA(b)).
\end{equation}
Both are independent of the representative, and $dF=0$. The integer
cochain is required to be closed; there is no restriction on its flux.
The curvature therefore remains a real, unbounded field.

Put $r_q=2\pi/q$. Each class has a unique representative with
$A(b)\in[0,r_q)$ at every link. The measure on the quotient is defined
by
\begin{equation}\label{eq:reg-configuration-integral}
 \int_{\operatorname{Conf}(\Gamma)}\Phi(A,z)\,\dd(A,z)
 =\sum_{\substack{z\in C^2(\Gamma;\Z)\\dz=0}}
       \int_{[0,r_q)^{3V}}\Phi(A,z)\,\dd A.
\end{equation}
This formula holds for nonnegative measurable functions and for
absolutely integrable complex functions. Any measurable fundamental
domain may be used instead. Indeed, the changes of representative
preserve Lebesgue measure in $A$ and counting measure in $z$.

We will repeatedly use the finite number
\begin{equation}\label{eq:reg-gaussian-sum}
 C_q(\varepsilon)=\sup_{y\in\R}\sum_{j\in\Z}
                  e^{-\varepsilon(y+r_qj)^2},\qquad \varepsilon>0.
\end{equation}
It satisfies
\[
 C_q(\varepsilon)\le2+\frac{q}{2\pi}\sqrt{\frac{\pi}{\varepsilon}}.
\]
To see this, retain the two terms nearest the origin and bound each
remaining term by $r_q^{-1}$ times the Gaussian integral over the
adjacent interval of length $r_q$ closer to the origin. Those intervals
have disjoint interiors. Write $C_q=C_q(1/2)$. Dropping the constraint
$dz=0$ in the flux sum gives
\begin{equation}\label{eq:reg-maxwell-integral}
 0<\int_{\operatorname{Conf}(\Gamma)}e^{-\|F\|^2/2}\,\dd(A,z)
 \le(r_qC_q)^{3V}.
\end{equation}
The lower bound follows already from $z=0$; for the upper bound, sum
each $z_\sigma$ separately at fixed $A$.

The periodic extensions of the curvatures lie in
\[
 \Omega=\{F\in C^2:dF=0\},
\]
with the topology of pointwise convergence on plaquettes. This is a
closed subspace of a countable product of copies of $\R$. Local
observables depend on finitely many plaquette values. For a finitely
supported real two-cochain $\beta$, the character
$\exp(i\langle\beta,F\rangle)$ is bounded and continuous on $\Omega$.

\subsection{Chambers and the fermion graph}\label{subsec:reg-chambers}

The lattice mirrors are
\begin{equation}\label{eq:reg-mirrors}
 \{x_\mu=c\},\quad c\in\tfrac12\Z,
 \qquad
 \{x_\mu\pm x_\nu=c\},\quad c\in\Z,\quad\mu<\nu.
\end{equation}
They form three coordinate families and six diagonal families. Let
$\mathcal L$ be the corresponding linear forms $\ell$ and write
$\ell(x)=n_\ell\cdot x$. The admissible positions are
$\operatorname{Adm}(x_\mu)=\frac12\Z$ and
$\operatorname{Adm}(x_\mu\pm x_\nu)=\Z$. The reflection in a mirror
$\Pi=(\ell,c)$ is
\begin{equation}\label{eq:reg-reflection}
 \theta_\Pi x=x-\frac{2(\ell(x)-c)}{|n_\ell|^2}n_\ell.
\end{equation}
Its upward normal is $n_\ell/|n_\ell|$. Coordinate mirrors at integer
positions are called site planes, and those at half-integer positions
are called link planes.

The group generated by these reflections is
\[
 W=\Z^3\rtimes O_{\mathrm{cub}},
 \qquad O_{\mathrm{cub}}=O(3)\cap\operatorname{GL}(3,\Z).
\]
Here $O_{\mathrm{cub}}$ consists of the $48$ signed permutation matrices.
A fundamental chamber is
\[
 C_0=\{0<x_3<x_2<x_1<\tfrac12\}.
\]
Its wall reflections are
\begin{align*}
 s_0(x)&=(x_1,x_2,-x_3),&
 s_1(x)&=(x_1,x_3,x_2),\\
 s_2(x)&=(x_2,x_1,x_3),&
 s_3(x)&=(1-x_1,x_2,x_3).
\end{align*}
Denote the walls by $\Pi_i$ and their inward unit normals by
\begin{equation}\label{eq:reg-wall-normals}
 n_0=e_3,\quad n_1=(e_2-e_3)/\sqrt2,
 \quad n_2=(e_1-e_2)/\sqrt2,\quad n_3=-e_1.
\end{equation}

Fix $v_0=(v_1,v_2,v_3)\in C_0$. The fermion sites and their neighbours
are
\begin{equation}\label{eq:reg-fermion-sites}
 \mathcal X=Wv_0,\qquad (gv_0)^{(i)}=gs_iv_0.
\end{equation}
The unordered bond $\{x,x^{(i)}\}$ has type $i$. A site has the unique
form $x=k+rv_0$, with $k\in\Z^3$ and $r\in O_{\mathrm{cub}}$; put
$\pi(x)=k$ and $y(x)=rv_0$. Thus $\pi(x)$ is its nearest photon site.

\begin{proposition}\label{prop:reg-geometry}
The chambers are precisely $gC_0$, $g\in W$, and $W$ acts on them
simply transitively. No fermion site lies on a mirror. Every bond
crosses exactly one mirror, its chamber wall, and crosses it
orthogonally at its midpoint. Consequently a bond crossing $\Pi$ is
$\{x,\theta_\Pi x\}$.

The projection $\pi$ is $W$-equivariant. Bonds of types $0,1,2$ have
the same photon site at both ends; a type-$3$ bond joins two nearest
photon sites. Exactly eight fermion bonds lie over each photon link.
If a mirror does not cut the link and its open segment lies on one
side, all eight bonds lie on that side.
The graph is connected and bipartite, and each site has one neighbour
of each type.
\end{proposition}

\begin{proof}
Coordinate reflections and coordinate transpositions generate
$O_{\mathrm{cub}}$, and
$\theta_{x_\mu,1/2}\theta_{x_\mu,0}=\tau_{e_\mu}$ gives the
translations. Conversely, \eqref{eq:reg-reflection} is an affine signed
permutation with integer translation, so the mirror group is $W$.

For a point on no mirror, subtract the unique nearest lattice point.
The resulting three coordinates are nonzero, have distinct absolute
values, and lie in $(-1/2,1/2)$. There is exactly one signed
permutation making them positive and decreasing. This places the
point in exactly one translate of $C_0$. These translates are disjoint
open convex sets, meet no mirror, and cover its complement. They are
therefore the chambers, with the asserted free transitive action.

For the four bonds from $v_0$, the midpoint is, respectively,
\[
 (v_1,v_2,0),\quad
 (v_1,(v_2+v_3)/2,(v_2+v_3)/2),\quad
 ((v_1+v_2)/2,(v_1+v_2)/2,v_3),\quad
 (1/2,v_2,v_3).
\]
Each belongs to the relative interior of its wall and to no other
mirror. Each half of the bond, except its midpoint, lies in one of
the two adjacent chambers. Transport by $g\in W$ proves the bond
assertion at every site.

In the coordinates $x=k+rv_0$, the first three neighbours are
$k+rs_iv_0$, whereas the fourth is
$k+re_1+rr_{e_1}v_0$, where $r_{e_1}$ changes the sign of the first
coordinate. This proves the statements about $\pi$ and its
equivariance. Over $(k,\mu)$ the lower ends are exactly the eight
sites for which $y_\mu=v_1$: the remaining two absolute values can
be ordered in two ways, with four choices of their signs. Their
type-$3$ neighbours are the upper ends of these eight bonds.

The bond over $b=(k,\mu)$ has the segment
$k+te_\mu+y_\perp$, $v_1\le t\le1-v_1$, where the other two
coordinates of $y_\perp$ have absolute values $v_2,v_3$. For a
diagonal height involving $x_\mu$, its varying part is
$t\pm y_\nu\in(0,1)$, the same interval between integers as on the
open photon link. For heights not involving $x_\mu$, the changes
are $y_\nu$ or $y_\nu\pm y_{\nu'}$, whose absolute values are
smaller than $1/2$ or $1$, respectively. Away from a mirror
containing $b$ these preserve its side. The coordinate height
$x_\mu$ gives the same side except at the unique link plane through
the midpoint, which cuts $b$. This proves the assertion about
photon links.

The generators $s_0,s_1,s_2$ generate the signed permutations, and
$s_3r_{e_1}=\tau_{e_1}$; conjugation gives all coordinate translations.
Thus the four wall reflections generate $W$, proving connectivity.
The two parts are the sites $gv_0$ with $\det r=1$ and with
$\det r=-1$, since each neighbour changes this determinant.
\end{proof}

The geometric location records reflection, while $(k,r)$ labels the
fermion sites by a photon site and a signed permutation. If
$d_i=\operatorname{dist}(v_0,\Pi_i)$, then
\[
 x^{(i)}-x=-2d_i\,rn_i,
 \qquad \frac{x^{(i)}-x}{|x^{(i)}-x|}=-rn_i.
\]
The directions, bond types and projected endpoints in these
coordinates do not depend on $v_0$. Hence neither the operator below
nor its determinant depends on $v_0$ after this identification.
\subsection{The chamber operator and the finite-volume law}
\label{subsec:reg-law}

Let $\sigma_\mu$ be the Pauli matrices and
\[
 \gamma_\mu=\diag(\sigma_\mu,-\sigma_\mu),\qquad
 \gamma_v=\sum_\mu v_\mu\gamma_\mu.
\]
They form a Hermitian Clifford triple. For a unit vector $n$, put
$P_n^\pm=(1\pm\gamma_n)/2$. Both projections have rank two, and
\begin{equation}\label{eq:reg-clifford-reflection}
 \gamma_n\gamma_v\gamma_n=-\gamma_{r_nv},\qquad
 r_nv=v-2(v\cdot n)n.
\end{equation}
This follows by multiplying
$\gamma_n\gamma_v+\gamma_v\gamma_n=2n\cdot v$ by $\gamma_n$.

For an oriented bond $(x,x')$ of type $i$ define
\begin{equation}\label{eq:reg-hop}
 U_A(x,x')=\exp(iqA(\pi(x),\pi(x'))),\qquad
 w(x,x')=\frac{w_i}{2}\left(\rho-\gamma_{(x'-x)/|x'-x|}\right).
\end{equation}
The chamber operator is
\begin{equation}\label{eq:reg-dirac}
 (\mathcal D_A\psi)(x)=M\psi(x)-\sum_{i=0}^3
                   w(x,x^{(i)})U_A(x,x^{(i)})\psi(x^{(i)}).
\end{equation}
Its lattice mass is
\begin{equation}\label{eq:reg-lattice-mass}
 m=M-\frac\rho2\sum_iw_i.
\end{equation}
A fermion boundary condition is a character
$\chi:\Gamma\to\{1,-1\}$. Write $\mathbb T=(\Gamma,\chi)$.
Its spinors satisfy $\psi(x+g)=\chi(g)\psi(x)$ for $g\in\Gamma$.
The restriction $\mathcal D_A^{\mathbb T}$ acts on a space of dimension $192V$,
with the inner product obtained by summing over representatives of
$\mathcal X/\Gamma$. For each site introduce generators
$\psi_\alpha(x),\bar\psi_\alpha(x)$, $1\le\alpha\le4$, obeying the same
boundary conditions. Their Grassmann algebra is denoted by
$\operatorname{Gr}_{\mathbb T}$. Its Berezin integral is normalized by
\begin{equation}\label{eq:reg-berezin}
 \int_{\mathbb T}\prod_{x\in\mathcal X/\Gamma}
                  \prod_{\alpha=1}^4\psi_\alpha(x)\bar\psi_\alpha(x)=1.
\end{equation}
Products of pairs are even, so their ordering is immaterial. With
$S_{\mathbb T}(A)=\sum_x\bar\psi(x)(\mathcal D_A\psi)(x)$, this normalization
gives $\int_{\mathbb T}e^{-S_{\mathbb T}(A)}=\det \mathcal D_A^{\mathbb T}$.

The unnormalized functional is
\begin{equation}\label{eq:reg-functional}
 Z_{\mathbb T}[G]=\int_{\operatorname{Conf}(\Gamma)}
      e^{-\|F\|^2/2}\int_{\mathbb T}e^{-S_{\mathbb T}(A)}G(A,z)
      \,\dd(A,z).
\end{equation}
For a scalar $G$ the inner integral is
$G(A,z)\det \mathcal D_A^{\mathbb T}$. For a Grassmann-valued integrand,
absolute convergence is understood coefficient by coefficient.
Bounded measurable coefficients suffice: $e^{-S_{\mathbb T}(A)}$
has finitely many coefficients, each a polynomial in the link phases,
and \eqref{eq:reg-maxwell-integral} supplies an integrable majorant.

\begin{proposition}\label{prop:reg-determinant-sign}
For every real $A$ and $M$,
\[
 \det \mathcal D_A^{\mathbb T}\ge0,
 \qquad
 \|\mathcal D_A^{\mathbb T}\|\le |M|+\frac{1+\rho}{2}\sum_iw_i.
\]
Moreover $\Re \mathcal D_A^{\mathbb T}=m+\frac\rho2\nabla_A^*\nabla_A$, where
\[
 (\nabla_A\psi)(x,i)=\sqrt{w_i/2}\,
           \bigl(\psi(x)-U_A(x,x^{(i)})\psi(x^{(i)})\bigr).
\]
The determinant is not identically zero as a function of the link
phases, for every real mass parameter $M$.
\end{proposition}

\begin{proof}
Write
\[
 H_A\psi(x)=\sum_iw_iU_A(x,x^{(i)})\psi(x^{(i)}),\qquad
 J_A\psi(x)=\sum_iw_i\gamma_{(x^{(i)}-x)/|x^{(i)}-x|}
                         U_A(x,x^{(i)})\psi(x^{(i)}).
\]
Reversing a bond conjugates its phase and reverses its direction.
Reindexing by the involution $x\mapsto x^{(i)}$ therefore gives
$H_A^*=H_A$ and $J_A^*=-J_A$. Hence
\[
 \mathcal D_A=M-\frac\rho2H_A+\frac12J_A,\qquad
 \nabla_A^*\nabla_A=\sum_iw_i-H_A.
\]
This proves the real-part identity. In the splitting
$\C^4=\C^2\oplus\C^2$, the two blocks of $H_A$ coincide, and the two
blocks of $J_A$ have opposite signs. It follows that
\begin{equation}\label{eq:reg-two-blocks}
 \mathcal D_A^{\mathbb T}=\diag(\check{\mathcal D}_A^{\mathbb T},
                                      (\check{\mathcal D}_A^{\mathbb T})^*),
 \qquad \det \mathcal D_A^{\mathbb T}=|\det\check{\mathcal D}_A^{\mathbb T}|^2.
\end{equation}
Each hopping matrix has norm $w_i(1+\rho)/2$, while the associated
translation with phase is unitary. Their sum gives the norm bound.

For the last assertion we view the determinant as a Laurent
polynomial $P_{\mathbb T}$ in the independent variables $u_b$.
Only type-$3$ bonds contain a variable, with $u_b$ at the lower end
and $u_b^{-1}$ at the upper end. Fix a coordinate direction $\mu$.
The coefficient of
\begin{equation}\label{eq:reg-extreme-monomial}
 \prod_{b\parallel e_\mu}u_b^{32}
\end{equation}
will be computed without using $M$.

Partition the sites modulo $\Gamma$ into six layers according to
the increasing values
\[
 y_\mu=-v_1,-v_2,-v_3,v_3,v_2,v_1.
\]
Each layer has $8V$ sites. Each site has one upward neighbour,
one downward neighbour, and two neighbours in its own layer. The
upward map is a permutation of the sites and visits successive layers
modulo six. Its bond types are
\begin{equation}\label{eq:reg-upward-types}
 (i_1,i_2,i_3,i_4,i_5,i_6)=(2,1,0,1,2,3).
\end{equation}
For layers $1,\ldots,5$ the upward phase is one; in layer six it is
$u_b$ for the link $(\pi(x),\mu)$. These facts follow by changing the
smallest sign or exchanging successive absolute values in the
neighbour formulas above. The sixth step changes $v_1$ to $-v_1$
and increases the photon coordinate by one.

There are eight lower and eight upper fermion ends over each photon
link, each carrying four matrix rows. An exponent $32$ at a link
therefore requires every lower row to take the upward type-$3$
block and every upper row to avoid the reverse block. Imposing this
at all links in direction $\mu$ forces every row in every layer to
take its upward block. Indeed, layer-six rows fill all columns in
layer one. A layer-one row cannot then use a column in its own layer,
and the reverse block is excluded by the exponent condition; it must
go upward to layer two. Inductively the columns in the current and
preceding layers are already occupied, so layers two through five
must also go upward. The selected terms use no other link variables.

The remaining freedom is a permutation of the four spin indices at
each site. The site permutation, repeated four times, has sign one.
Summing over the spin permutations therefore gives the product of
the $4\times4$ determinants of the selected hopping blocks. The
boundary signs have fourth power one, and
\[
 \det w(x,x^\uparrow)
   =(w_{i_k}/2)^4(1-\rho^2)^2
\]
at a site in layer $k$. Consequently the coefficient of
\eqref{eq:reg-extreme-monomial} is $C_*^V$, where
\begin{equation}\label{eq:reg-extreme-coefficient}
 C_*=(w_0/2)^{32}(w_1/2)^{64}(w_2/2)^{64}(w_3/2)^{32}
                                           (1-\rho^2)^{96}>0.
\end{equation}
This coefficient is independent of $M$. In particular the determinant
polynomial, and its restriction to the phase torus, are nonzero.
\end{proof}

The real-part estimate alone would require control of the sign of
$m$ to prove invertibility. The coefficient computation instead
proves that the determinant has positive integral at every real $M$,
which is all that is needed to define the law.

\begin{proposition}\label{prop:reg-partition}
The partition function satisfies $0<Z_{\mathbb T}[1]<\infty$.
More precisely,
\begin{equation}\label{eq:reg-partition-bounds}
 r_q^{3V}e^{-6r_q^2V}C_*^V
 \le Z_{\mathbb T}[1]
 \le(r_qC_q)^{3V}\left(|M|+\frac{1+\rho}{2}\sum_iw_i\right)^{192V}.
\end{equation}
Thus $\E_{\mathbb T}[G]=Z_{\mathbb T}[G]/Z_{\mathbb T}[1]$ defines a
probability law on scalar observables. Its curvature law
$\mu_{\mathbb T}$ is the law of the periodic extension $\widetilde F$
on $\Omega$.
\end{proposition}

\begin{proof}
For a nonnegative Laurent polynomial $P$ on the phase torus, every
Fourier coefficient satisfies
\[
 |\widehat P(\lambda)|
 =\left|\int P(u)\overline{u^\lambda}\,\dd u\right|
 \le\int P(u)\,\dd u,
\]
where Haar measure is normalized. By the preceding proposition and
\eqref{eq:reg-extreme-coefficient},
\[
 \int_{[0,r_q)^{3V}}\det \mathcal D_A^{\mathbb T}\,\dd A
 \ge r_q^{3V}C_*^V.
\]
On this box $|(dA)_\sigma|\le2r_q$ and
$\|dA\|^2\le12Vr_q^2$. The $z=0$ term in
\eqref{eq:reg-functional} gives the lower bound in
\eqref{eq:reg-partition-bounds}. For the upper bound use
$|\det \mathcal D_A^{\mathbb T}|\le\|\mathcal D_A^{\mathbb T}\|^{192V}$ and
\eqref{eq:reg-maxwell-integral}. The density
$e^{-\|F\|^2/2}\det \mathcal D_A^{\mathbb T}$ is nonnegative. After division
by its strictly positive integral it is a probability density, whose
image under $(A,z)\mapsto\widetilde F$ is the asserted law.
\end{proof}

